\documentclass[english]{article}
\usepackage{listings}
\usepackage{todonotes}
\setlength{\parindent}{0in}
\title{Journal of Object Technology Submission JOT-2013-014: Cover Letter}
\sloppy
\begin{document}
\maketitle
We thank the reviewers for their valuable feedback on this paper. We
revisited our submission to address the issues highlighted in the
reviews. In this letter, we describe our actions to meet the
requirements for resubmission.  For better support of our ideas, we
included some more figures and sections, re-organized the section
order, and extended explanations in several sections.  In the
following we address specific questions of the reviewers. Comments of
the reviewers are set \textit{italic}. Our response is set
\textbf{bold}.

\subsection*{Reviewer 1}
\textit{The system has 1.5 slowdown even if constraints are not used.
  Could this be eliminated using an extension of PyPy's type inference
  to detect variables that are not involved in constraints?}\\
\textbf{It cannot be determined statically which variables are not
  involved in constraints. However, the JIT learns which variables are
  constrained and which are ordinary object-oriented variables. In
  compiled machine code, the check is removed where possible and no
  overhead remains (in the absence of constraints, the generated
  assembler is exactly the same as for the unmodified VM). In the
  presented benchmarks, any overhead is in code that the JIT has not
  compiled to machine code. The maximum slowdown of 1.5 is for the
  \emph{binary trees} benchmark, where little code is compiled. We
  have clarified this in the section
  describing the benchmark results.}\\

\textit{The system has not been evaluated on solvers that find entire
  data structures such as arrays or strings. The constraint solving is
  only applied to numerical domains in
  the two example benchmarks presented.}\\
\textbf{We have yet to find good benchmarks in these areas. However,
  we provide a solver on arrays with numeric elements, that can
  express constraints over the elements, as well as the size of the
  array. We have included this in the section ``Adding New
  Solvers''. Furthermore, in the appendix, the electrical circuits
  example uses constraints on arrays to constrain the currents in a
  list of connected leads to sum to 0. We noted this in the
  explanation, but have added a line reference, now}\\

% \textit{First, there are only two constraint solving benchmarks
%   presented [sic] It is my opinion that one of the main factors
%   holding back progress in constraint solving support in languages is
%   the lack of good challenge benchmarks, so the paper could do more to
%   contribute to this.}\\
% \todo{figure out what to answer}\textbf{}\\

\textit{the presentation and the design mix the programming model
  and its implementation...}\\
\textbf{We have restructured the paper. We now present the theory,
  then introduce our semantic changes first in an implementation
  independent manner and then give examples, and finally, we describe
  only what we consider the most important parts of the
  implementation. We have also made clearer which features we consider
  part of the language, and which features we consider are solver
  specific.}\\

\textit{I would have also liked to see a performance comparison with
  other recent systems, such as Kaplan system presented in POPL 2012
  paper ``Constraints as Control'' [...] The performance comparison is
  interesting because Kaplan uses Scala, which is statically typed,
  and interoperates with Z3 reasonably
  efficiently.}\\
\textbf{We ran a few puzzle solving benchmarks to compare our
  performance with Kaplan, the Turtle language, and Prolog. We're
  faster than Kaplan in those, generally slower than Prolog (which
  isn't very surprising, given that the benchmarks use no imperative
  features, and are really better suited for Prolog.) However, we have
  found that different versions of the Z3 library can vary wildly in
  their performance, and sometimes just changing the order in which
  constraints are declared can have a big impact. We have tried to
  change the order for each language to exhibit its best performance.}\\

\textit{The authors gloss over the questions of finite precision of
  numbers used by some solvers.  Z3 does not support floating points,
  only real numbers, whose semantics is very different than floating
  points.}\\
\textbf{This was initially ignored, because the Z3 C-API returns
  results as double precision floating point numbers. We now mention
  this problem for Z3. Currently, programmers have to be aware that
  there may be precision loss here. Cassowary uses floats, and uses a
  hard-coded $\epsilon$ to determine if two floats should be
  considered equal. We have not, in practice, had problems with this,
  but we agree that the question remains how much the programmer has
  to be aware of precision issues. We will investigate this further,
  but conclude for now that the programmer has to be aware that
  rounding issues may occur when using floats. Z3 also supports
  integers, and Babelsberg/R represents Ruby Fixnums as Z3 integers if
  possible.}\\

\textit{As far as I can see, the presented system does not support
  enumeration of solutions. It is possible to find all
  solutions when the set is finite?}\\
\textbf{No, enumerating solutions is not currently available. We have
  made it clearer in the comparison with BackTalk, Ilog, and Kaplan in
  the Related Work section, as well as in the theoretical background
  section, that this is not among the primary goals of Babelsberg. We
  think that enumerating solutions will be useful for debugging and
  plan to implement it primarily for that use case. We also plan to
  provide Kodkod as a finite domain solver, and enumerating solutions
  will be useful here as well. For now, however, in Babelsberg we are
  interested in \emph{a} solution, so if there
  are multiple, the solvers are expected to select one arbitrarily.}\\

\textit{I was also wondering if identity constraints are useful for
  finding data structures?}\\
\textbf{If we supported enumerating solutions, possibly, yes. This
  would be similar BackTalk or Squander, where a constraint can be
  satisfied by searching for appropriate objects on the heap or
  creating objects and testing them for suitability. In our case, it
  is useful for modeling data structures where the difference between
  "we need two values to be equal" and "two values need to refer to
  the \emph{same} object" is important. This is an issue, for example,
  in a Doubly Linked List, where \lstinline{self.previous.next} should
  be \emph{identical} to \lstinline{self}, not just equal. We have
  added a reference to the Kaleidoscope paper on object identity and
  provide a short example.}\\

\textit{Could the authors provide a reference to the DeltaBlue library
  and examples where they used it?}\\
\textbf{We now introduce solvers that are interesting for our
  discussion in the introduction, among them DeltaBlue. We have added
  an example of using it in the section ``Solving Constraints''.}\\

\subsection*{Reviewer 2}

There were a number of comments on the choice of words, structure of
sentences, and grammatical errors. Additionally, there were requests
for line references in source code listings, as well as suggestions
for explanations, shortening some parts, and references to papers to
include. We have applied these improvements, and do not mention them
in our answers below.

\textit{The biggest problem with the paper is that the language is
  \_only\_ explained through examples. The semantics, not even syntax,
  of the extensions are explained in any general level, so it is
  really not possible to know what Babelsberg is.  There is a section
  on implementation, but that, on the other hand, is too detailed, and
  vague at the same time. I don't believe that it is possible to infer
  the semantics, or even syntax, from that description either.}\\
\textbf{The syntactic extension is now explained as being optional at
  the beginning of the design section and in the overview. The syntax
  of the always construct is the same as for a Ruby method call. We
  now refer a description in the appendix.\\
  We have re-structured the paper extensively to introduce the theory
  first, then explain our two semantic extensions first in an
  implementation independent way, then with an example. The
  implementation section has been shortened to focus on what we
  consider the important details.}\\

\textit{You have not defined what "soft constraints" are.}\\
\textbf{We introduce these terms earlier now, with a section on
  relevant constraint theory.}\\

\textit{Why in the example are both farenheit and celcius initialized?
  Would the constraint not take care of initializing one from the
  other?}\\
\textbf{Indeed it does, we have changed the code to avoid this
  confusion.}\\

\textit{``Here these are just accessors''\\
  - Here seems to refer to "constraints" but that doesn't seem to make
  sense.  Constraints are not accessors?}\\
\textbf{``Here'' was supposed to refer to the example, ``these'' to
  the methods ``centigrade'' and ``fahrenheit'', but the sentence
  structure was misleading. We have clarified it.}\\

\textit{What is the syntax for specifying which solver to use?}\\
\textbf{It is simply an additional argument to the \lstinline{always}
  method. We have added example code for this in the ``Solving
  Constraints'' section, and describe it in the appendix}\\

\textit{The reference to "co-operating constraint solvers" is
  confusing. [...] this is how I thought your system already works?
  Based on a type of the variables, the constraint solver to apply is
  selected(?)}\\
\textbf{The system already supports different solvers in the same
  program and selects which solver to use based on the type of the
  variable. However, different solvers cannot collaborate to find a
  value for the same variable -- each variable is always determined by
  one particular solver. The cooperating solvers architecture allows
  multiple solvers to collaboratively determine the value of the same
  variable. The reference has been removed, because this is part of
  the future work and only confuses.}\\

\textit{"or imperatively"\\
  - does this mean "assigning to the variable directly"\\
  - if so, would it not then violate its constraints?\\}\\
\textbf{We have added some explanation. The constraints are checked on
  assignment, and the ``read-only'' constraint is updated. A
  ``read-only'' annotation does not ``freeze'' a variable, it just
  prevents the solvers from changing the variable value to satisfy
  constraints. This is also explained in more detail in the
  section on relevant theory.}\\

\textit{"If multiple variables that have constraints on them are
  assigned using multi-assignment, all values are assigned before the
  solver is
  triggered."\\
  - Could this lead to conflicting (unsatisfiable) constraints if
  those variables are involved with more than one constraint?}\\
\textbf{Yes. As with normal assignment, if because of these assignment
  the constraints cannot be satisfied, the assignments are not applied
  and an exception is raised (which the programmer can handle.) We
  have clarified this in the text. The reason for not applying the
  assignments in that case is, that the programmer can continue after
  an exception and the system is still in a valid state.}

\textit{You use priority ``high''. Does it have an exact meanign? (or
  any meaning?) [...] What does priority ``low'' mean?  (In relation
  to ``high'', and is there ``must'' or ``even lower'' or
  ``middle''..., and is
  there an ordering between the ``low''s?)}\\
\textbf{We have added a detailed explanation of the relevant constrain
  theory. The theory allows an arbitrary number of preferences, but
  prior work has found that only few priorities tend to be used in
  stylized ways. The symbolic names are just for comprehension. There
  is no ordering between different ``low''s in Babelsberg, so the
  solver is free to trade them off arbitrarily.}\\

\textit{In 3.7, the example you give about incremental constraints
  seems incomplete.  It is difficiult for me to see how the
  TemperatureConverter system could benefit from any incremental
  solving? Can you provide some more details?}\\
\textbf{The idea is that the input to the slider, instead of an
  enumerator, may be a graphical slider. We have added this to the
  text.}\\

\textit{is the semantics of your ``edit'' language construct that it
  invokes an incremental solver if one is available? }\\
\textbf{We have added the explanation that ``edit'' is a method
  supplied by a particular solver. This means that the developer has
  to know about such a method for a particular solver, in our case
  Cassowary, and then has to be explicit about using it.}\\

\textit{The section about implementation is not useful.}\\
\textbf{We have shortened the implementation section. Many of the
  details are omitted. Other details, which we consider crucial to
  understanding our implementation, have been expanded and changed to
  be more comprehensible, and to focus on the relevant issues.}\\

\textit{Please first give the full syntax of an edit constraint
  specification, then describe its semantics, and only then try to to
  explain the implementaiton (if you think it is necessary).}\\
\textbf{We have split the explanation of edit constraints into theory,
  design in Babelsberg, and the relevant implementation details in
  Babelsberg/R. We hope this clarifies that edit constraints have no
  special syntax, and are just methods defined by the libraries that
  support them.}\\

\textit{``Considering that Babelsberg/R is aimed at imperative
  programmers, we expect that most of the time the vm will not be
  solving constraints in tight loops, but running mostly imperative
  code intermingled with
  constraint re-satisfaction.''\\
  - Are you sure about this assumption?}\\
\textbf{We removed this sentence. We did not mean to imply constraints
  would not be used. However, in the same way that expensive
  operations are mostly not used in tight loops in object-oriented
  programs (e.g., most call sites are mono-morphic, reflective
  features are rarley used -- assumptions which are the reason why
  JITs work so well), we believe that object-constraint programs
  likely do something besides repeatedly solving constraints. In the
  video streaming example, besides solving constraints for every
  frame's quality, each frame is also encoded and written to a
  stream.}\\

\subsection*{Reviewer 3}

Where the more detailed comments suggested specific text changes, we
have improved the relevant sections. These are not mentioned below.

% That said, I found it hard to identify what is the novel research
% contribution.
%   TIM TODO: put the contributions in more contrast to kaleidoscope
%   TODO: Object Constraint Programming section, added a sentence to contrast our goals against Kaleidoscope

\textit{While the authors have picked Ruby, and they probably need to
  pick a specific language, can't this be done to other OO languages
  or dynamic languages. [...] Why is Ruby chosen
  anyway? Why not other OO languages? Why not Java even?}\\
\textbf{Yes. We have mentioned a JavaScript implementation of
  Babelsberg as future work. We now mention this earlier, to make
  clear these ideas are not tied to Ruby. The reason for choosing Ruby
  was that one of the authors was already familiar with the
  VM.}\\

\textit{The authors should compare the OO features of Babelsberg/R
  with Ilog Solver even though the CIP focus does make this paper
  different.}\\
\textbf{Ilog Solver, like Kaplan or Backtalk, has a different focus,
  which we now expand upon in the ``Related Work'' section.}\\

\textit{Notice that Ilog Solver does not need to have different
  execution modes like what is done here because of the use of
  overloading.}\\
\textbf{Ilog Solver uses overloading, for example, to be able to do
  arithmetic with numbers and \emph{decision variables}. However, Ilog
  Solver cannot create a constraint expression from expressions
  involving only ordinary numbers. Instead, the programmer has to
  explicitly create decision variables to build constraints. These
  decision variables can then be used (through overloading) in
  conjunction with normal numbers. However, Babelsberg does not use
  execution modes in that case either (the VM just reads constrained
  variables differently and
  dispatches to the type of their value.)\\
  The execution modes are used to avoid having to explicitly create
  \emph{decision variable}-like objects. This allows constraints in
  Babelsberg to use domain specific abstractions available on the
  constrained objects, without having to
  add further overloads to a \emph{decision variable} class.}\\

\textit{Very few constraint solvers are designed for CIP as it
  requires specific treatment of soft constraints. [...] Its not clear
  how those solvers, e.g., Z3 will fit well. It would be better to
  have a clearer treatment of solver semantics combined with language
  semantics.}\\
\textbf{While soft constraints are useful for many application
  domains, CIP does not necessarily require them. The section on
  ``Preferential Constraints'' now explains this, and also explains
  that priorities can be interpreted by the solver. In our
  implementation, Z3 will raise an exception if a priority is
  passed.}\\

\textit{Having a formalization which can encompass several solvers and
  their different features, [...] can then be clearly mapped to both
  syntax, semantics and implementation in Babelsberg/R.}\\
\textbf{The section ``Solving Constraints'' now elaborates that solver
  specifics (such as numeric precision, how they select solutions, or
  whether they support incremental solving or soft constraints) are
  not considered part of the language. Instead, programmers choose
  solvers as they choose among any multiple alternative
  libraries -- by weighing their properties.\\
  The current contribution is not a unification of different solvers'
  features, but the unification of declaring and solving constraints
  with object-oriented methods and abstractions. The ability to use
  ordinary methods in constraints allows programmers to use libraries
  and domain specific abstractions in constraints without having to
  express their meaning in both an imperative and a declarative
  way.}\\

\textit{Note that in S2.1, the Z3 example is quite different from the
  other examples.}\\
\textbf{This example is not Babelsberg code, but ordinary Ruby code
  using the Z3 library. We have improved the section to clarify
  this.}\\

\textit{What happens if an object is no longer live (and hence can be
  deleted) but the constraints interact with another constraint in
  another live object. Normally the variables in the dead object can
  be garbage collected and also the data structures but this may not
  be correct with constraints.}\\
\textbf{A description of GC issues is now included in the
  implementation section. A short explanation: normal variables are
  garbage collected normally. Constrained variables are referenced
  from the variable storage locations in the VM, so if those go out of
  scope, they can be collected, and their last value remains in the
  solver specific data structures. However, variables \emph{can}
  additionally be referenced from the block closures that created the
  constraints, so additional references may exist. For objects that
  have no solver representation, the blocks are kept alive (which
  would be most complex data structures). This means that the
  variables are kept alive, too}\\

% \textit{The contribution of the CIP part of the paper is mentioned
%   only briefly. Is it mainly the implementation?}\\
% \textbf{The sections regarding the simpler semantic model and the
%   interaction with the solvers have}\\

\textit{It is unclear to me in the evaluation section that the
  imperative version is behaviorally equivalent to the
  object-constraint version. Shouldn't the constraint version be more
  general than the imperative version?  If so the comparison may not
  be totally fair for the constraint
  version.}\\
\textbf{Yes, its not totally fair, but in practice we may not care for
  the more general solution, we just care about making it more concise
  or expressive. If we assume the different versions came about as
  part of a refactoring, then we think the comparison is important. We
  have added this argument to the text.}\\

\textit{Are the Edit Constraints tied to the Cassowary solver or can
  one use other solvers?}\\
\textbf{We have added an explanation that ``edit'' is a method in the
  Cassowary library. This means that the developer has to know about
  this method and then has to be explicit about using it.}\\

% \textit{It would be good to have a single example which can be run
%   with different solvers and also compared with imperative code
%   (perhaps the imperative code might be big so may need to be in
%   appendix).}\\
% \todo{?}\textbf{}\\

\end{document}

%%% Local Variables:
%%% mode: latex
%%% End:
